\rok{Новембар 2023.}{A}

Први практични тест из Алгоритама и структура података

\zadatak{1}{Insertion Sort}
Написати \texttt{InsertionSort} алгоритам за сортирање целобројног низа дужине \textit{n}.

\textbf{Додатне напомене за решавање задатка:}
\begin{itemize}
    \item Улаз је несортиран низ целих бројева.
    \item Излаз је сортиран низ целих бројева.
    \item У шаблону се налази класа \texttt{RandomArrayGenerator} коју треба искористити за генерисање низа случајних бројева.
    \item Креирати класу \texttt{InsertionSort} и у оквиру ње написати имплементацију за методу:
    \begin{Verbatim}[fontsize=\footnotesize, numbers=left, xleftmargin=10mm]
public static void Sort(long[] arr)
    \end{Verbatim}
\end{itemize}



\zadatak{2}{Да ли је неки број палиндром?}
Креирање функције која ће проверавати да ли је неки број палиндром. Палиндром је број који остаје идентичан уколико му се цифре окрену. Алгоритам је неопходно написати са линеарном временском комплексношћу O(n).

\textbf{Додатни захтеви за решавање задатка:}
\begin{itemize}
    \item Алгоритам писати у оквиру класе \texttt{DoublyLinkedList}.
    \item Захтеван потпис методе:
    \begin{Verbatim}[fontsize=\footnotesize, numbers=left, xleftmargin=10mm]
public bool IsDoublyLinkedListPalindrome()
    \end{Verbatim}
\end{itemize}

\textbf{Примери:}
\begin{itemize}
    \item \textbf{Пример 1:} \\
          \textbf{Улаз:} $2 \leftrightarrow 3 \leftrightarrow 3 \leftrightarrow 2$ \\
          \textbf{Излаз:} true
    \item \textbf{Пример 2:} \\
          \textbf{Улаз:} $2 \leftrightarrow 3 \leftrightarrow 2$ \\
          \textbf{Излаз:} true
    \item \textbf{Пример 3:} \\
          \textbf{Улаз:} $2 \leftrightarrow 3 \leftrightarrow 1 \leftrightarrow 1 \leftrightarrow 3 \leftrightarrow 2$ \\
          \textbf{Излаз:} false
\end{itemize}



\zadatak{3}{Брисање узастопних парова из низа бројева}
Креирати алгоритам који из низа целих бројева брише све узастопне парове вредности. Узастопни парови подразумевају вредности које се налазе на суседним позицијама у низу. Елементи из низа се бришу искључиво уколико се у проласку кроз низ два иста броја нађу један поред другог.

\textbf{Додатни захтеви за решавање задатка:}
\begin{itemize}
    \item Задатак решавати помоћу структуре стека (класа \texttt{Stack}) која је дата у оквиру шаблона задатка.
    \item Захтеван потпис методе:
    \begin{Verbatim}[fontsize=\footnotesize, numbers=left, xleftmargin=10mm]
public static void ShortenArray(int[] array)
    \end{Verbatim}
\end{itemize}

\textbf{Примери:}
\begin{itemize}
    \item \textbf{Пример 1:}
    \begin{itemize}
        \item \textbf{Улаз:} [1, 2, 2, 3, 3, 1]
        \item \textbf{Кораци брисања елемената у низу:}
        \begin{itemize}
            \item[] {[1 2 2 3 3 1]}
            \item[] {[2 2 3 3]}
            \item[] {[3 3]}
            \item[] {[]}
        \end{itemize}
        \item \textbf{Излаз:} All pairs are found!
    \end{itemize}
    \item \textbf{Пример 2:}
    \begin{itemize}
        \item \textbf{Улаз:} [2, 2, 1, 4, 3, 1]
        \item \textbf{Кораци брисања елемената у низу:}
        \begin{itemize}
            \item[] {[2 2 1 4 3 1]}
            \item[] {[1 4 3 1]}
        \end{itemize}
        \item \textbf{Излаз:} Remained numbers are: 1 4 3 1
    \end{itemize}
\end{itemize}



\sablon{AISP2023_Test1_GrupaA}{2023/Novembar/GrupaA/AISP2023_Test1_GrupaA.linq}